\notechapter{N-20260314}{A Cookbook for Multivariable Topology, Limits, and Differential Computations}


\section{Classifying points relative to a set}

For a set $G\subseteq\mathbb R^n$, classify a point $P$ by inspecting small balls around it.  If some ball lies inside $G$, $P$ is interior.  If some ball is disjoint from $G$, $P$ is exterior.  If every ball meets both $G$ and $G^c$, $P$ is a boundary point.  This is the fastest way to understand open, closed, and neither-open-nor-closed examples.

\section{Limits}

To prove a multivariable limit, dominate the expression by a function of $r=\|x-a\|$ going to zero.  To disprove it, use two paths.  Polar coordinates are useful near the origin because they separate radial size from angular dependence:
\[
  x=r\cos\theta,\qquad y=r\sin\theta.
\]
If a candidate expression becomes $r^\alpha g(\theta)$ with bounded $g$ and $\alpha>0$, the limit is zero.

\section{Differentiability and tangent objects}

The standard hierarchy is
\[
  C^1\Rightarrow\text{differentiable}\Rightarrow\text{continuous}\Rightarrow\text{directional limits behave well in value}.
\]
The reverse implications usually fail.  Tangent lines, normal lines, and tangent planes are all first-order approximations.  For $F(x,y,z)=0$, the normal vector is $\nabla F$; for $z=f(x,y)$, the tangent plane comes from the differential $dz=f_xdx+f_ydy$.

\section{Expanded synthesis and advanced context}

A cookbook is useful only if it remembers why the recipes work.  The recipes here are all forms of one idea: replace a complicated local object by its simplest stable approximation, either a ball for topology, a radial estimate for limits, or a linear map for differentiability.



\paragraph{Expanded synthesis.}
For this chapter, the elementary material should be read as a study of what remains invariant when coordinates change. The earlier sections (`Classifying points relative to a set`, `Limits`, `Differentiability and tangent objects`, `Higher viewpoint`, `Expanded cookbook with failure modes`) compute with arrays, bases, pivots, determinants, or eigenvectors; the deeper object is a linear map together with the spaces it acts on. A matrix is a representation of that map after bases have been chosen. This is why formulas such as
\[
  A' = P^{-1}AP,\qquad \widetilde A = T^{-1}AS
\]
are not cosmetic: they distinguish an intrinsic morphism from a coordinate description.

The structural hierarchy is as follows. Kernels record what a map forgets, images record what it reaches, cokernels record obstructions to solvability, and quotients record information after an equivalence has been imposed. Determinants act on the top exterior power and therefore measure oriented volume; traces are invariant under cyclic rearrangement and become characters in representation theory; adjoints depend on an inner product and lead to self-adjoint spectral theory.

A useful way to return to the computations is to ask, for every formula, which invariant it is measuring. Row reduction measures rank and kernel; diagonalization measures decomposition into invariant subspaces; Gram--Schmidt measures orthogonal projection; positive definiteness measures ellipsoid geometry; tensor notation measures how components transform. The elementary methods are therefore the finite-dimensional laboratory in which duality, quotienting, functoriality, and spectral decomposition can be seen clearly.


\section{Expanded cookbook with failure modes}

A cookbook is useful only when each recipe includes its failure mode.

For limits, path testing can disprove a limit but cannot prove it unless all paths are covered by an estimate.  For continuity, partial continuity in each variable is necessary but insufficient.  For differentiability, existence of all directional derivatives is also insufficient unless they assemble into a linear approximation.

\section{Practical limit templates}

Near the origin, convert to polar coordinates when the expression is homogeneous or nearly homogeneous.  If
\[
  f(r\cos\theta,r\sin\theta)=r^\alpha g(\theta)
\]
with $\alpha>0$ and bounded $g$, then the limit is zero.  If $g(\theta)$ remains in the limit with no decaying factor, the limit depends on direction.

For rational expressions, compare total degrees.  If the numerator has strictly higher degree than the denominator and the denominator is comparable to a power of $r$, the limit is often zero.  But cancellations and anisotropic denominators require care.

\section{Tangent and normal templates}

For $z=f(x,y)$, use the graph formula.  For $F(x,y,z)=0$, use the gradient.  For a space curve given as the intersection of $F=0$ and $G=0$, the tangent direction is
\[
  \nabla F\times \nabla G.
\]
This follows because the tangent vector must be perpendicular to both normals.
